\SourceLeaf{152}
II. Quum $b^m$ sit residuum absolute minimum ipsius $-b^{m-1}$ secundum mod.\ $a^m$, maior quam $\frac12 a^m$ non erit (art.\ 4).

III. Quia $a^m a^{m+1}=D+b^m b^m$, atque $a^{m+1}$ non $<a^m$, $a^m a^m$ non erit $>D+b^m b^m$, et quum $b^m$ non $>\frac12 a^m$, $a^m a^m$ non erit $>D+\frac14 a^m a^m$ et $\frac34 a^m a^m$ non $>D$, tandemque $a^m$ \emph{non} $>\surd{\frac43 D}$.

\emph{Exempl.} Proposita sit forma $(304,217,155)$, cuius determinans $=-31$. Hic invenitur progressio formarum:
\[
(304,217,155),\ (155,-62,25),\ (25,12,7),\ (7,2,5),\ (5,-2,7).
\]
Ultima est quaesita. --- Eodem modo formae $(121,49,20)$, cuius determinans $=-19$, aequivalentes inveniuntur: $(20,-9,5)$, $(5,-1,4)$, $(4,1,5)$: quare $(4,1,5)$ erit forma quaesita.

Tales formas $(A,B,C)$, quarum determinans est negativus et in quibus $A$ nec maior quam $\surd{\frac43 D},C$, nec minor quam $2B$, \emph{formas reductas} vocabimus. Quare cuivis formae determinantis negativi, forma reducta proprie aequivalens inveniri poterit.

\Needspace{8\baselineskip}
\Article{172}
\textsc{Problema.} \emph{Invenire conditiones, sub quibus duae formae reductae non identicae, eiusdem determinantis $-D$, $(a,b,c)$, $(a',b',c')$ proprie aequivalentes esse possint.}

\emph{Solutio.} Supponamus, id quod licet, $a'$ esse non $>a$, formamque $axx+2bxy+cyy$ transire in $a'x'x'+2b'x'y'+c'y'y'$ per substitutionem propriam $x=\alpha x'+\beta y'$, $y=\gamma x'+\delta y'$. Tum habebuntur aequationes
\[
\begin{aligned}
a\alpha\alpha+2b\alpha\gamma+c\gamma\gamma&=a'&\quad\mathinner{.\,.\,.\,.\,.\,.\,.\,.}\quad&[1]\\
a\alpha\beta+b(\alpha\delta+\beta\gamma)+c\gamma\delta&=b'&\quad\mathinner{.\,.\,.\,.\,.\,.\,.}\quad&[2]\\
\alpha\delta-\beta\gamma&=1&\quad\mathinner{.\,.\,.\,.\,.\,.\,.\,.\,.\,.\,.}\quad&[3]
\end{aligned}
\]
Ex [1] sequitur $aa'=(a\alpha+b\gamma)^2+D\gamma\gamma$; quare $aa'$ erit positivus; et quum $ac=D+bb$, $a'c'=D+b'b'$, etiam $ac$, $a'c'$ positivi erunt: quare $a,a',c,c'$ omnes eadem signa habebunt. Sed tum $a$ tum $a'$ non $>\surd{\frac43 D}$, adeoque $aa'$ non $>\surd{\frac43 D}$; quare multo minus $D\gamma\gamma\ (=aa'-(a\alpha+b\gamma)^2)$ maior quam $\frac43 D$ esse poterit. Hinc $\gamma$ erit aut $=0$, aut $=\pm1$.

I. Si $\gamma=0$, ex [3] sequitur esse aut $\alpha=1.\,\delta=1$, aut $\alpha=-1$, $\delta=-1$.
\SourceLeaf{153}
In utroque casu fit ex [1] $a'=a$, et ex [2] $b'-b=\pm\beta a$. Sed $b$ non $>\frac12 a$, et $b'$ non $>\frac12 a'$, proin etiam non $>\frac12 a$. Quare aequatio $b'-b=\pm\beta a$ consistere nequit, nisi fuerit

\emph{aut} $b=b'$; unde sequeretur $c'=\frac{b'b'+D}{a'}=\frac{bb+D}{a}=c$, quare formae $(a,b,c)$, $(a',b',c')$ identicae essent contra hyp.

\emph{aut} $b=-b'=\pm\frac12 a$. In hoc etiam casu erit $c'=c$ formaque $(a',b',c')$ erit $(a,-b,c)$ \emph{i.\ e.} formae $(a,b,c)$ opposita. Simul patet formas has esse ancipites propter $2b=\pm a$.

II. Si $\gamma=\pm1$, fit ex [1] $a\alpha\alpha+c-a'=\pm2b\alpha$. Sed $c$ non minor quam $a$, adeoque non minor quam $a'$: hinc $a\alpha\alpha+c-a'$ sive $2b\alpha$ certo non minor quam $a\alpha\alpha$. Quare quum $2b$ non sit maior quam $a$, erit $\alpha$ non minor quam $\alpha\alpha$; unde necessario aut $\alpha=0$, aut $=\pm1$.

1) Si $\alpha=0$, fit ex [1] $a'=c$, et quoniam $a$ neque maior quam $c$, neque minor quam $a'$, erit necessario $a'=a=c$. Porro ex [3] fit $\beta\gamma=-1$, unde ex [2] $b+b'=\pm\delta c=\pm\delta a$. Hinc simili modo ut in (I) sequitur esse

\emph{aut} $b=b'$, in quo casu formae $(a,b,c)$, $(a',b',c')$ forent identicae, contra hyp.

\emph{aut} $b=-b'$, in quo casu formae $(a,b,c)$, $(a',b',c')$ erunt oppositae.

2) Si $\alpha=\pm1$, ex [1] sequitur $\pm2b=a+c-a'$. Quare quum neque $a$. neque $c<a'$, erit $2b$ non $<a$, et non $<c$. Sed $2b$ etiam non $>a$. neque $>c$, unde necessario $\pm2b=a=c$, et hinc ex aequ.\ $\pm2b=a+c-a'$, etiam $=a'$. Fit igitur ex [2]
\[
b'=a(\alpha\beta+\gamma\delta)+b(\alpha\delta+\beta\gamma)
\]
sive, propter $\alpha\delta-\beta\gamma=1$,
\[
b'-b=a(\alpha\beta+\gamma\delta)+2b\beta\gamma=a(\alpha\beta+\gamma\delta\pm\beta\gamma)
\]
quare necessario, ut ante

\emph{aut} $b=b'$, unde formae $(a,b,c)$, $(a',b',c')$ identicae, contra hyp.

\emph{aut} $b=-b'$, adeoque formae illae oppositae. Simul in hoc casu propter $a=\pm2b$, formae erunt ancipites.

Ex his omnibus colligitur, formas $(a,b,c)$, $(a',b',c')$ proprie aequivalentes esse non posse nisi fuerint oppositae, simulque \emph{aut} ancipites, \emph{aut} $a=c=a'=c'$.
\SourceLeaf{154}
In hisce casibus formas $(a,b,c)$, $(a',b',c')$ proprie aequivalere, vel a priori facile praevideri potuit; si enim formae sunt oppositae, improprie, et si insuper ancipites, etiam proprie aequivalentes esse debent; si vero $a=c$, forma $(\frac{D+(a-b)^2}{a},a-b,a)$ formae $(a,b,c)$ contigua et proin aequivalens erit; sed propter $D+bb=ac=aa$ fit $\frac{D+(a-b)^2}{a}=2a-2b$, forma vero $(2a-2b,a-b,a)$ est anceps; quare $(a,b,c)$ oppositae suae etiam proprie aequivalebit.

Aeque facile iam diiudicari potest, quando duae formae reductae $(a,b,c)$, $(a',b',c')$ non oppositae improprie aequivalentes esse possint. Erunt enim impr.\ aequivalentes, si $(a,b,c)$, $(a',-b',c')$, quae non identicae erunt, proprie sunt aequivalentes, et contra. Hinc patet, conditionem, sub qua \emph{illae} improprie sint aequivalentes, esse, ut sint identicae, insuperque aut ancipites aut $a=c$. --- Formae vero reductae quae neque identicae sunt neque oppositae, neque proprie neque improprie aequivalentes esse possunt.

\Needspace{8\baselineskip}
\Article{173}
\textsc{Problema.} \emph{Propositis duabus formis eiusdem determinantis negativi, $F$ et $F'$, investigare utrum sint aequivalentes.}

\emph{Solutio.} Quaerantur duae formae reductae $f$, $f'$ formis $F$, $F'$ resp.\ proprie aequivalentes: si formae $f$, $f'$ sunt proprie, vel improprie vel utroque modo aequivalentes, etiam $F$, $F'$ erunt; si vero $f$, $f'$ nullo modo aequivalentes sunt, etiam $F$, $F'$ non erunt.

Ex art.\ praec.\ dari possunt quatuor casus:

1) Si $f$, $f'$ neque identicae neque oppositae, $F$, $F'$ nullo modo aequivalentes erunt.

2) Si $f$, $f'$ sunt \emph{primo} vel identicae vel oppositae, et \emph{secundo} vel ancipites, vel terminos suos extremos aequales habent: $F$, $F'$ tum proprie, tum improprie aequivalentes erunt.

3) Si $f$, $f'$ sunt identicae, neque vero ancipites neque terminos extremos aequales habent: $F,F'$ proprie tantum aequivalebunt.

4) Si $f$, $f'$ sunt oppositae, neque vero ancipites, neque terminos extremos aequales habent: $F,F'$ improprie tantum aequivalentes erunt.

\emph{Ex.} Formis $(41,35,30)$, $(7,18,47)$ quarum determinans $=-5$, reductae $(1,0,5)$, $(2,1,3)$ aequivalentes inveniuntur, quare illae nullo modo aequivalentes erunt. --- Formis vero $(23,38,63)$, $(15,20,27)$ aequivalet eadem
\SourceLeaf{155}
reducta $(2,1,3)$, quae quum simul sit anceps, formae $(23,38,63)$, $(15,20,27)$ tum proprie tum improprie aequivalebunt. --- Formis $(37,53,78)$, $(53,73,102)$, aequivalent reductae $(9,2,9)$, $(9,-2,9)$ quae quum sint oppositae, ipsarumque termini extremi aequales: formae propositae tam proprie quam improprie erunt aequivalentes.

\Needspace{8\baselineskip}
\Article{174}
Multitudo omnium formarum reductarum, determinantem datum $-D$ habentium, semper est finita, et, respectu numeri $D$, satis modica: formae hae ipsae vero duplici modo inveniri possunt. Designemus formas reductas determinantis $-D$ indefinite per $(a,b,c)$, ubi itaque omnes valores ipsorum $a,b,c$ determinari debent.

\emph{Methodus prima.} Accipiantur pro $a$ omnes numeri. tum positivi tum negativi non maiores quam $\surd{\frac43 D}$, quorum residuum quadraticum $-D$, et pro singulis $a$, fiat $b$ successive aequalis omnibus valoribus expr.\ $\surd{-D}\ (\mathrm{mod.}\ a)$, non maioribus quam $\frac12 a$, tum positive tum negative acceptis; $c$ vero pro singulis valoribus determinatis ipsorum $a,b$, ponatur $=\frac{D+bb}{a}$. Si quae formae hoc modo oriuntur in quibus $c<a$, hae erunt reiiciendae, reliquae autem manifesto erunt reductae.

\emph{Methodus secunda.} Accipiantur pro $b$ omnes numeri, tum positivi tum negativi, non maiores quam $\frac12\surd{\frac43 D}$, sive $\surd{\frac13 D}$; pro singulis $b$ resolvatur $bb+D$ omnibus quibus fieri potest modis in binos factores (etiam signorum diversitatis ratione habita) ambos ipso $2b$ non minores ponaturque alter factor, et quidem, quando factores sunt inaequales, minor $=a$, alter $=c$. Quum $a$ non erit $>\surd{\frac43 D}$, omnes formae quae hoc modo prodeunt, manifesto erunt reductae. --- Denique patet, nullam formam reductam dari posse quae non per utramque methodum inveniatur.

\emph{Ex.} Sit $D=85$. Hic limes valorum ipsius $a$ est $\surd{\frac{340}{3}}$ qui iacet inter $10$ et $11$. Numeri vero inter $1$ et $10$ (incl.) quorum residuum $-85$, sunt $1,2,5,10$. Unde habentur formae duodecim:
\[
\begin{gathered}
(1,0,85),\ (2,1,43),\ (2,-1,43),\ (5,0,17),\\
(10,5,11),\ (10.\,-5,11);\ (-1,0,-85),\ (-2,1,-43),\\
(-2,-1,-43).\ (-5,0,-17),\ (-10,5,-11,\ (-10,-5,-11).
\end{gathered}
\]
Per methodum alteram limes valorum ipsius $b$ habetur $\surd{\frac{85}{3}}$, qui situs est inter $5$ et $6$. Pro $b=0$, prodeunt formae
\SourceLeaf{156}
\[
(1,0,85),\ (-1,0,-85),\ (5,0,17),\ (-5,0,-17),
\]
pro $b=\pm1$ hae $(2,\pm1,43)$, $(-2,\pm1,-43)$.

Pro $b=\pm2$ nullae habentur, quia $89$ in duos factores, qui ambo non $<4$, resolvi nequit. Idem valet de $\pm3$, $\pm4$. Tandem pro $b=\pm5$, proveniunt
\[
(10,\pm5,11),\ (-10,\pm5.\,-11).
\]

\Needspace{8\baselineskip}
\Article{175}
Si ex omnibus formis reductis determinantis dati, formarum binarum, quae, licet non identicae, tamen proprie sunt aequivalentes, alterutra reiicitur: formae remanentes hac insigni proprietate erunt praeditae, ut, quaevis forma eiusdem determinantis alicui ex ipsis proprie sit aequivalens, et quidem unicae tantum (alias enim inter ipsas aliquae proprie aequivalentes forent). Unde patet, \emph{omnes formas eiusdem determinantis in totidem classes distribui posse quot formae remanserint}, referendo scilicet formas eidem reductae proprie aequivalentes in eandem classem. Ita pro $D=85$, remanent formae
\[
\begin{gathered}
(1,0,85),\ (2,1,43),\ (5,0,17),\ (10,5,11),\\
(-1,0.\,-84),\ (-2,1,-43).\ (-5,0,-17),\ (-10,5,-11)
\end{gathered}
\]
quare omnes formae determinantis $-85$ in octo classes distribui poterunt, prout formae primae, aut secundae etc.\ proprie aequivalent. Perspicuum vero est, formas in eadem classe locatas proprie aequivalentes fore, formas ex diversis classibus proprie aequivalentes esse non posse. Sed hoc argumentum de classificatione formarum infra multo fusius exsequemur. Hic unicam observationem adiicimus. Iam supra ostendimus, si determinans formae $(a,b,c)$ fuerit negativus $=-D$, $a$ et $c$ eadem signa habere (quia scilicet $ac=bb+D$ adeoque positivus); eadem ratione facile perspicitur, si formae $(a,b,c)$, $(a',b',c')$ sint aequivalentes, omnes $a,c,a',c'$ eadem signa habituros. Si enim prior in posteriorem per substitut.\ $x=\alpha x'+\beta y'$, $y=\gamma x'+\delta y'$ transit: erit $a\alpha\alpha+2b\alpha\gamma+c\gamma\gamma=a'$, hinc $aa'=(a\alpha+b\beta)^2+D\gamma\gamma$, adeoque certo non negativus; quoniam vero neque $a$ neque $a'=0$ esse potest, erit $aa'$ positivus et proin signa ipsorum $a,a'$ eadem.

Hinc manifestum est; formas quarum termini exteri sint positivi, ab iis quarum termini exteri sint negativi, prorsus esse separatas, sufficitque ex formis reductis eas tantum considerare quae terminos suos exteros positivos habent, nam
\SourceLeaf{157}
reliquae totidem sunt multitudine et ex illis oriuntur, tribuendo terminis exteris signa opposita; idemque valet de formis ex reductis reiiciendis et remanentibus.

\Needspace{8\baselineskip}
\Article{176}
Ecce itaque pro determinantibus quibusdam negativis tabulam formarum, secundum quas omnes reliquae eiusdem determinantis in classes distingui possunt; apponimus autem, ad annotat.\ art.\ praec., semissem tantum, scilicet eas quarum termini exteri positivi.
\[
\begin{array}{r|l}
\multicolumn{1}{r}{D}&\\
1&(1,0,1).\\
2&(1,0,2).\\
3&(1,0,3),\ (2,1,2).\\
4&(1,0,4),\ (2,0,2).\\
5&(1,0,5),\ (2,1,3).\\
6&(1,0,6),\ (2,0,3).\\
7&(1,0,7),\ (2,1,4).\\
8&(1,0,8),\ (2,0,4),\ (3,1,3).\\
9&(1,0,9),\ (2,1,5),\ (3.\,0,3).\\
10&(1,0,10),\ (2,0,5).\\
11&(1,0,11),\ (2,1,6),\ (3,1,4),\ (3,-1,4).\\
12&(1,0,12),\ (2,0,6),\ (3,0,4),\ (4,2,4).
\end{array}
\]
Superfluum foret hanc tabulam hic ulterius continuare. quippe quam infra multo aptius disponere docebimus.

Patet itaque, quamvis formam determinantis $-1$, formae $xx+yy$ proprie aequivalere, si ipsius termini exteri sint positivi, vel huic $-xx-yy$. si sint negativi; quamvis formam determinantis $-2$, cuius termini exteri positivi, formae $xx+2yy$ etc.; quamvis formam determinantis $-11$, cuius termini exteri positivi, alicui ex his $xx+11yy$, $2xx+2xy+6yy$, $3xx+2xy+4yy$, $3xx-2xy+4yy$ etc.

\Needspace{8\baselineskip}
\Article{177}
\textsc{Problema.} \emph{Habetur series formarum, quarum quaevis praecedenti a parte posteriori contigua: desideratur transformatio aliqua propria primae in formam quamcunque seriei.}

\SourceLeaf{158}
\emph{Solutio.} Sint formae $(a,b,a')=F$; $(a',b',a'')=F'$; $(a'',b'',a''')=F''$; $(a''',b''',a'''')=F'''$ etc.\ Designentur $\frac{b+b'}{a'}$, $\frac{b'+b''}{a''}$, $\frac{b''+b'''}{a'''}$ etc.\ respective per $h',h'',h'''$ etc.\ Sint indeterminatae formarum $F,F',F''$ etc.\ $x,y$; $x',y'$; $x'',y''$ etc.\ Ponatur $F$ transmutari
\[
\begin{array}{ll@{\qquad}l@{\qquad}l}
\text{in }F'&\text{positis}&x=\alpha'x'+\beta'y',&y=\gamma'x'+\delta'y'\\
F''&\ldots&x=\alpha''x''+\beta''y'',&y=\gamma''x''+\delta''y''\\
F'''&\ldots&x=\alpha'''x'''+\beta'''y''',&y=\gamma'''x'''+\delta'''y'''\\
&&\text{etc.}&
\end{array}
\]
{\small\setlength{\arraycolsep}{3pt}
\[
\begin{array}{l@{\enspace}l@{\enspace}l}
\text{Tum quia }F\text{ transit in }F'\text{ positis}&x=-y',&y=x'+h'y'\\
\hfill F'\text{ in }F''\text{ positis}&x'=-y'',&y'=x''+h''y''\\
\hfill F''\text{ in }F'''\text{ positis}&x''=-y''',&y''=x'''+h'''y'''\text{ etc. (art.\ 160)}
\end{array}
\]
}
facile eruetur sequens algorithmus (art.\ 159):
\[
\begin{array}{ll@{\qquad}ll}
\alpha'=0&\beta'=-1&\gamma'=1&\delta'=h'\\
\alpha''=\beta'&\beta''=h''\beta'-\alpha'&\gamma''=\delta'&\delta''=h''\delta'-\gamma'\\
\alpha'''=\beta''&\beta'''=h'''\beta''-\alpha''&\gamma'''=\delta''&\delta'''=h'''\delta''-\gamma''\\
\alpha''''=\beta'''&\beta''''=h''''\beta'''-\alpha'''&\gamma''''=\delta'''&\delta''''=h''''\delta'''-\gamma'''\\
&\text{etc.}&&
\end{array}
\]
sive
\[
\begin{array}{ll@{\qquad}ll}
\alpha'=0&\beta'=-1&\gamma'=1&\delta'=h'\\
\alpha''=\beta'&\beta''=h''\beta'&\gamma''=\delta'&\delta''=h''\delta'-1\\
\alpha'''=\beta''&\beta'''=h'''\beta''-\beta'&\gamma'''=\delta''&\delta'''=h'''\delta''-\delta'\\
\alpha''''=\beta'''&\beta''''=h''''\beta'''-\beta''&\gamma''''=\delta'''&\delta''''=h''''\delta'''-\delta''\\
&\text{etc.}&&
\end{array}
\]

Omnes has transformationes esse proprias tum ex ipsarum formatione tum ex art.\ 159 nullo negotio deduci potest.

Algorithmus hic perquam simplex et ad calculum expeditus, algorithmo in art.\ 27 exposito est analogus, ad quem etiam reduci potest\HistoricalNote{*)}{Erit scilicet in signis art.\ 27, $\beta^n=\pm[-h'',h''',-h''''\ldots\pm h^n]$ ubi signa ambigue posita, esse debent $-\,-$; $-\,+$; $+\,-$; $+\,+$; prout $n$ formae $4k+0$; $1$; $2$; $3$, et
\[
\delta^n=\pm[h',-h'',h'''\ldotp\ldotp\ldotp\ldotp\ldotp\pm h^n]
\]
ubi signa ambigua esse debent $+\,-$; $+\,+$; $-\,-$; $-\,+$, prout $n$ formae $4k+0$; $1$; $2$; $3$. Sed hoc, quod quivis facile ipse confirmare poterit, fusius exsequi, nobis brevitas non permittit.}. Ceterum solutio
\SourceLeaf{159}
haec ad formas determinantis negativi non est restricta, sed ad omnes casus patet, si modo nullus numerorum $a',a'',a'''$ etc.\ $=0$.

\Needspace{8\baselineskip}
\Article{178}
\textsc{Problema.} \emph{Propositis duabus formis $F$, $f$, eiusdem determinantis negativi, proprie aequivalentibus: invenire transformationem aliquam propriam alterius in alteram.}

\emph{Sol.} Supponamus formam $F$ esse $(A,B,A')$ et per methodum art.\ 171 inventam esse progressionem formarum $(A',B',A'')$, $(A'',B'',A''')$ etc.\ usque ad $(A^m,B^m,A^{m+1})$, quae sit reducta: similiterque $f$ esse $(a,b,a')$ et per eandem methodum inventam seriem $(a',b',a'')$, $(a'',b'',a''')$ usque ad $(a^n,b^n,a^{n+1})$, quae sit reducta. Tum duo casus locum habere possunt.

I. Si formae $(A^m,B^m,A^{m+1})$, $(a^n,b^n,a^{n+1})$ sunt aut identicae, aut oppositae simulque ancipites. Tum formae $(A^{m-1},B^{m-1},A^m)$, $(a^n,-b^{n-1},a^{n-1})$ erunt contiguae (designante $A^{m-1}$ terminum progressionis $A,A',A''\ldots A^m$ penultimum, similiaque $B^{m-1}$, $a^{n-1}$, $b^{n-1}$). Nam $A^m=a^n$, $B^{m-1}\equiv-B^m\ (\mathrm{mod.}\ A^m)$, $b^{n-1}\equiv-b^n\ (\mathrm{mod.}\ a^n\text{ sive }A^m)$, unde $B^{m-1}-b^{n-1}\equiv b^n-B^m$ adeoque $\equiv0$, si formae $(A^m,B^m,A^{m+1})$, $(a^n,b^n,a^{n+1})$ sunt identicae, et $\equiv2b^n$ adeoque $\equiv0$, si sunt oppositae et ancipites. Quare in progressione formarum
\[
\begin{gathered}
(A,B,A'),\ (A',B',A'')\ldotp\ldotp\ldotp\ldotp(A^{m-1},B^{m-1},A^m),\\
(a^n,-b^{n-1},a^{n-1}),\ (a^{n-1},-b^{n-2},a^{n-2})\ldots(a',-b,a),\ (a,b,a')
\end{gathered}
\]
quaevis forma praecedenti contigua erit, adeoque per art.\ praec.\ transformatio propria primae $F$ in ultimam $f$ inveniri poterit.

II. Si formae $(A^m,B^m,A^{m+1})$, $(a^n,b^n,a^{n+1})$ non identicae, sed oppositae simulque $A^m=A^{m+1}=a^n=a^{n+1}$. Tum progressio formarum
\[
\begin{gathered}
(A,B,A'),\ (A',B',A'')\ldotp\ldotp\ldotp\ldotp(A^m,B^m,A^{m+1}),\\
(a^n,-b^{n-1},a^{n-1}),\ (a^{n-1},-b^{n-2},a^{n-2})\ldotp\ldotp\ldotp\ldotp(a',-b,a),\ (a,b,a')
\end{gathered}
\]
eadem proprietate erit praedita. Nam $A^{m+1}=a^n$, et $B^m-b^{n-1}=-(b^n+b^{n-1})$ per $a^n$ divisibilis. Unde per art.\ praec.\ invenietur transformatio propria formae primae $F$ in ultimam $f$.

\SourceLeaf{160}
\emph{Ex.} Ita pro formis $(23.\,38,63)$, $(15,20,27)$ habetur progressio
\[
\begin{gathered}
(23,38,63).\ (63,25,10),\ (10,5,3),\ (3,1.\,2),\\
(2.\,-7,27),\ (27,-20,15),\ (15,20,27)
\end{gathered}
\]
quare
\[
h'=1,\quad h''=3,\quad h'''=2,\quad h''''=-3,\quad h'''''=-1,\quad h''''''=0.
\]
Hinc deducitur transformatio formae $23xx+76xy+63yy$ in $15tt+40tu+27uu$ haec: $x=-13t-18u$, $y=8t+11u$.

Ex solutione hac nullo negotio sequitur solutio problematis: \emph{Si formae $F,f$ improprie sunt aequivalentes, invenire transformationem impropriam formae $F$ in $f$.} Sit enim $f=att+2btu+a'uu$ eritque forma opposita $app-2bpq+a'qq$ formae $F$ proprie aequivalens. Quaeratur transformatio propria formae $F$ in illam, $x=\alpha p+\beta q$, $y=\gamma p+\delta q$, patetque $F$ transire in $f$ positis $x=\alpha t-\beta u$, $y=\gamma t-\delta u$, hancque transformationem fore impropriam.

Quodsi igitur formae $F,f$ tam proprie quam improprie sunt aequivalentes: inveniri poterit tam transformatio propria aliqua quam impropria.

\Needspace{8\baselineskip}
\Article{179}
\textsc{Problema.} \emph{Si formae $F$, $f$ sunt aequivalentes: invenire omnes transformationes formae $F$ in $f$.}

\emph{Sol.} Si formae $F$, $f$ unico tantum modo sunt aequivalentes \emph{i.\ e.} proprie tantum vel improprie tantum: quaeratur per art.\ praec.\ transformatio una formae $F$ in $f$, patetque alias quam quae huic sint similes dari non posse. Si vero formae $F$, $f$ tam proprie quam improprie aequivalent, quaerantur duae transformationes, altera propria, altera impropria. Iam sit forma $F=(A,B,C)$, $BB-AC=-D$, numerorumque $A,2B,C$ divisor communis maximus $=m$. Tum ex art.\ 162 patet, in priori casu omnes transformationes formae $F$ in $f$ ex una transformatione, in posteriori omnes proprias ex propria omnesque improprias ex impropria deduci posse, si modo omnes solutiones aequationis $tt+Duu=mm$ habeantur. His igitur inventis problema erit solutum.

Habetur autem $D=AC-BB$, $4D=4AC-4BB$, quare $\frac{4D}{mm}=4\frac{AC}{mm}-(\frac{2B}{m})^2$ erit integer. Iam si

1) $\frac{4D}{mm}>4$, erit $D>mm$: quare in $tt+Duu=mm$, $u$ necessario debebit esse $=0$, adeoque $t$ alios valores quam $+m$, et $-m$ habere nequit. Hinc
\SourceLeaf{161}
si $F$, $f$ unico tantum modo aequivalentes sunt et transformatio aliqua
\[
x=\alpha x'+\beta y',\quad y=\gamma x'+\delta y'
\]
praeter hanc ipsam quae prodit ex $t=m$ (art.\ 162), et hanc
\[
x=-\alpha x'-\beta y',\quad y=-\gamma x'-\delta y'
\]
aliae locum habere non possunt. Si vero $F$, $f$ tum proprie tum improprie aequivalent, atque propria aliqua transformatio habetur
\[
x=\alpha x'+\beta y',\quad y=\gamma x'+\delta y'
\]
impropriaque
\[
x=\alpha'x'+\beta'y',\quad y=\gamma'x'+\delta'y'
\]
praeter illam (ex $t=m$) et hancce
\[
x=-\alpha x'-\beta y',\ y=-\gamma x'-\delta y'\ (\text{ex }t=-m)
\]
alia propria non dabitur; similiterque nulla impropria praeter
\[
\begin{gathered}
x=\alpha'x'+\beta'y',\quad y=\gamma'x'+\delta'y';\\
\text{et }x=-\alpha'x'-\beta'y',\quad y'=-\gamma'x'-\delta'y'
\end{gathered}
\]

2) Si $\frac{4D}{mm}=4$, sive $D=mm$, aequatio $tt+Duu=mm$ quatuor solutiones admittet $t,u;=m,0;\ -m,0;\ 0,1;\ 0,-1$. Hinc si $F$, $f$ unico tantum modo aequivalentes et transformatio aliqua
\[
x=\alpha x'+\beta y',\quad y=\gamma x'+\delta y'
\]
\emph{quatuor} omnino transformationes dabuntur,
\[
x=\pm\alpha x'\pm\beta y',\quad y=\pm\gamma x'\pm\delta y'
\]
\[
\begin{aligned}
x&=\mp\frac{\alpha B+\gamma C}{m}x'\mp\frac{\beta B+\delta C}{m}y',\\
y&=\pm\frac{\alpha A+\gamma B}{m}x'\pm\frac{\beta A+\delta B}{m}y'
\end{aligned}
\]

Si vero $F$, $f$ duobus modis aequivalent, sive praeter transformationem illam datam alia ipsi dissimilis habetur: haec quoque suppeditabit quatuor illis dissimiles, ita ut \emph{octo} transformationes habeantur. --- Ceterum facile demonstrari potest in hoc casu $F$, $f$ semper revera duobus modis aequivalere. Nam quum $D=mm=AC-BB$, $m$ etiam ipsum $B$ metietur. Formae $(\frac{A}{m},\frac{B}{m},\frac{C}{m})$ determinans erit $=-1$, quare formae $(1,0,1)$ vel huic $(-1,0,-1)$ erit aequivalens,
